\section{Discussion and future work}\label{sec:discussion}

The theorem package is intentionally narrow in scope and correspondingly strong in closure: a closed package on an audited shell-stable class for a canonical hybrid Cantor object, rather than a broad theorem about all context-dependent Cantor systems, a shell-general theorem, or a theorem of non-SFT breadth.

\subsection{What is closed in the present paper}

The paper closes four theoremlets on the audited shell:
\begin{enumerate}
    \item a continuous full-loop lawfulness theoremlet,
    \item a cocycle pressure closure theoremlet,
    \item a strict theory extension theoremlet, and
    \item a conditional pressure disintegration theoremlet.
\end{enumerate}
Taken together, these theoremlets show that the audited shell carries a lawful continuous substrate, that the cocycle-level pressure object is closed there, that the completion-based packaged-object theory is a strict extension of the cocycle theory, and that this extension has a nontrivial thermodynamic consequence on the same shell-stable class.

\subsection{What the paper does not claim}

The scope restrictions are part of the mathematical content and should be read as theorem-level boundaries, not merely as stylistic caution. In particular, the paper does \emph{not} claim:
\begin{itemize}
    \item a broader-class theorem beyond the audited shell,
    \item a shell-general theorem,
    \item an external or non-SFT breadth theorem beyond what is actually closed here,
    \item a direct stratumwise root-separation theorem,
    \item a packaging-induced broader theorem class claim,
    \item or an exact KL-based closure-deficit theorem.
\end{itemize}
The support-only diagnostics retained in the appendix and internal package do not enlarge these claims. They accompany the theorem package, but they do not substitute for it and they do not widen its scope.

\subsection{Why audited-shell scope is still mathematically meaningful}

One reaction to the shell restriction is that the paper proves only a local result around a specially chosen regime. The shell restriction, however, is what allows the theorem object to remain coherent while preserving the full six-mechanism dynamics. A broader-class statement that silently dropped nondegenerate lens competition, packaging competition, budget activity, or timescale adaptation would describe a different and weaker object.

\subsection{Reserve and stretch routes}

Several mathematically meaningful routes remain outside the closed paper core.

\paragraph{Packaging completion as a broader class route.}
The packaging-completion object is real and mathematically nondegenerate, but the present paper does not claim that it yields a broader theorem class than the cocycle route on the audited shell. That possibility remains a reserve route.

\paragraph{Hybrid-object broadening beyond the audited shell.}
The paper does not claim that the canonical hybrid object already supports a shell-general theory. Whether the extension and disintegration theoremlets can be widened to a larger shell or a more canonical class remains open.

\paragraph{Wider parameter-shell results.}
The present result does not prove that the audited-shell class can be replaced by a substantially wider parameter shell. That remains future work.

\paragraph{External/non-SFT breadth.}
The paper does not claim that the present theorem package already establishes a non-SFT breadth result in the strong external sense. That ambition remains explicitly outside the current theorem package.

\paragraph{Alternative thermodynamic consequence routes.}
The direct stratumwise route was investigated and rejected as the wrong thermodynamic object for the present paper. The accepted route is conditional pressure disintegration. Whether stronger package-conditioned or measure-theoretic consequences can be obtained remains open.

\subsection{Future work}

There are four natural future directions.

\begin{enumerate}
    \item \textbf{Broader shell stability.}
    One can ask whether the audited shell can be widened while preserving all four theoremlets simultaneously.

    \item \textbf{A sharper theorem of non-definability.}
    The strict-extension theoremlet is closed on the audited shell, but a more intrinsic and possibly more canonical non-definability argument would strengthen the overall package.

    \item \textbf{Further thermodynamic consequences.}
    The paper closes conditional pressure disintegration, but stronger measure-, regularity-, or root-level consequences may exist on top of the same canonical object.

    \item \textbf{Broader object classes.}
    The most ambitious direction is to enlarge the theorem package beyond the present audited shell without collapsing the six-mechanism object structure that makes the result interesting in the first place.
\end{enumerate}

\subsection{Final perspective}

Rather than the broadest possible theorem class, the contribution is a closed theorem object on a lawful shell: a base cocycle theory, a completion-based extension, and a thermodynamic consequence of that extension. The natural direction for future work is to widen the class while preserving the same object-level architecture, not to trade this object for a broader but weaker statement.
